\documentclass[12pt]{article}

\usepackage[left=1in, right=1in, top=1in, bottom=1in]{geometry}
\usepackage{amsmath, amssymb, amsthm, mathtools, epigraph, bussproofs}

\DeclarePairedDelimiter{\norm}{\lVert}{\rVert}
\usepackage[T2A]{fontenc}
\usepackage[utf8]{inputenc}
\usepackage[russian]{babel}

\newtheorem{theorem}{Теорема}
\newtheorem{lemma}{Лемма}
\newtheorem{definition}{Определение}
\newtheorem{example}{Пример}
\newtheorem{corollary}{Следствие}
\newtheorem{algorithm}{Алгоритм}
\newtheorem{remark}{Замечание}
\newtheorem{notation}{Обозначение}

\title{Математическая логика}
\date{\today}


\setlength\epigraphwidth{.8\textwidth}
\setlength\epigraphrule{0pt}

\large
\begin{document}
\maketitle
\section{Язык высказываний}

Давайте перед началом вспомним, что такое формальный язык и контекстно-свободная грамматика.

Символы --- это значки, которые мы не определяем. Например, $a$, $b$, $c$ --- это символы.
Слова --- это последовательность символов ($abc$ --- это слово). Если мы пишем слова последовательно, то мы конкатенируем эти слова и получаем новое слово (Если $w = ab, v = cd $, то $ wv = abcd$). Мы бронием значок $\epsilon$ для обозначения пустого слова.

Алфавит --- это множество букв. Язык --- это множество слов.

Если $L$ --- язык, то $L^* = \bigcup_{i=0}^\infty L^i$

\begin{definition}
  Контекстно-свободная грамматика $G$ --- это $(V,T,P,S)$, где
  \begin{itemize}
  \item $V$ - мн-во нетерминальных символов,
  \item $T$ - мн-во терминальных символов,
  \item $P$ - мн-во продукций, конкретно, элементов $A \longrightarrow \alpha$, где $A \in V, \alpha \in (V \cup T)^*$
  \item $S \in V$ -  начальный символ

  \end{itemize}
  (Про то, как он порождает язык, я не буду говорить)))) )
\end{definition}

\begin{definition}
  Пусть $\mathcal{P}$ - произвольный алфавит, который мы будем называть множество пропозиональных переменных.

  Тогда $L$ - язык высказываний, если это язык в алфавите $\mathcal{P} \sqcup \{\land, \lor, \lnot, \rightarrow, ), (, \bot, \top\}$ и задаётся следующей грамматикой:
  \begin{align*}
    S & \longrightarrow p \text{, где } p \in \mathcal{P}\\
    & \longrightarrow \bot\\
    & \longrightarrow \top\\
    & \longrightarrow (\lnot S)\\
    & \longrightarrow (S \lor S)\\
    & \longrightarrow (S \land S)\\
    & \longrightarrow (S \to S)\\
  \end{align*}
\end{definition}

Договоримся, что у нас такой приоритет операций
$\lnot\land\lor\to$
, начиная, конечно, с первой операции.

Договоримся, что $\lor$ и $\land$ - левоассоциативные операция.
Насчёт $\to$: математики любят, чтобы они было левоассоциативными
, а программисты --- правоассоциативные. (Та же ассоциативность, что стрелочном типе в система типов HM).



\section{Классическая семантика}
%% \epigraph{\itshape Begin at the beginning, the King said gravely, ``and go on till you come to the end: then stop.''}{---Lewis Carroll, \textit{Alice in Wonderland}}
Приготовьтесь, здесь будет много определений.
  Здесь $\mathbb{B}_0 = \{0, 1\}$

\begin{definition}
  $\mathcal{I}: \mathcal{P} \rightarrow \mathbb{B}_0$ --- интерпретация (Напомним, что $\mathcal{P}$ --- это множество пропозициональных переменных)
\end{definition}

\begin{definition}
  $[\![ \phi ]\!]_{\mathcal{I}}$ - семантическая скобка, которая выдает значение из $\mathbb{B}_0$, где $\phi$ - формула
  Индуктивно определим:
  \begin{itemize}
  \item $[\![ p ]\!]_{\mathcal{I}} = \mathcal{I}(p)$, $p \in \mathcal{P}$
  \item $[\![ \bot ]\!]_{\mathcal{I}} = 0$
  \item $[\![ \top ]\!]_{\mathcal{I}} = 1$
  \item $[\![ \lnot \phi ]\!]_{\mathcal{I}} = 1 - [\![ \phi ]\!]_{\mathcal{I}}$
  \item $[\![ \phi \lor \psi ]\!]_{\mathcal{I}} = max([\![ \phi ]\!]_{\mathcal{I}}, [\![ \psi ]\!]_{\mathcal{I}})$
  \item $[\![ \phi \land \psi ]\!]_{\mathcal{I}} = min([\![ \phi ]\!]_{\mathcal{I}}, [\![ \psi ]\!]_{\mathcal{I}})$
  \item $[\![ \phi \to \psi ]\!]_{\mathcal{I}} = [\![ \lnot \phi \lor \psi ]\!]_{\mathcal{I}}$
  \end{itemize}
\end{definition}

\begin{example}
  Пусть $\mathcal{I}$ --- такая интерпретация, что $\mathcal{I}(p) = 1, \mathcal{I}(q) = 0$
  
  Тогда $[\![ (p \to q) \lor (q \to (q \land \lnot q)) ]\!]_{\mathcal{I}} = 1$
\end{example}

\begin{definition}
  \leavevmode  
  \begin{enumerate}
  \item $\phi$ - тавтология, если $\forall \ \mathcal{I}$ --- интерпретация, $[\![ \phi ]\!]_{\mathcal{I}} = 1$
  \item $\phi$ - невыполнима (UNSAT), если $\forall \ \mathcal{I}$ --- интерпретация, $[\![ \phi ]\!]_{\mathcal{I}} = 0$
  \item $\phi$ - выполнима (SAT), если $\exists \ \mathcal{I}$ --- интерпретация, $[\![ \phi ]\!]_{\mathcal{I}} = 1$

  \end{enumerate}
\end{definition}
\begin{example}
  В классической логике (которую мы рассматриваем) формула $a \lor \lnot a$ ---  тавтология
\end{example}

\begin{definition}[Логическое следствие]
  Если для всякой интерпретации $\mathcal{I}$, такое что $[\![ \phi ]\!]_{\mathcal{I}} = 1$, выполнено $[\![ \psi ]\!]_{\mathcal{I}} = 1$, то пишем $\phi \models \psi$.
\end{definition}

\begin{definition}
  $\mathcal{T}$ --- теория, если это множество формул.
\end{definition}

\begin{definition}
  \leavevmode
  \begin{enumerate}
  \item Теория $\mathcal{T}$ выполнима (SAT), если $\exists \ \mathcal{I}$ --- интерпретация, $\forall \ \phi \in \mathcal{T}, [\![ \phi ]\!]_{\mathcal{I}} = 1$.
  \item Теория $\mathcal{T}$ невыполнима (UNSAT), если $\forall \ \mathcal{I}$ --- интерпретация, $\exists \ \phi \in \mathcal{T}, [\![ \phi ]\!]_{\mathcal{I}} = 0$.
  \end{enumerate}
\end{definition}

\begin{definition}
  Пусть $\mathcal{T}$ --- теория, $\phi$ --- формула.
  Мы пишем $\mathcal{T} \models \phi$, если $\forall \ \mathcal{I}$ --- интерпретация, при которой $\mathcal{I}$ --- SAT, $[\![ \psi ]\!]_{\mathcal{I}} = 1$
\end{definition}

\begin{definition}[Модель теории]
  Пусть $\mathcal{T}$ --- теория, $\mathcal{I}$ --- интерпретация.

  
  $\mathcal{I}$ - модель теории $\mathcal{T}$, если $\forall \phi \in \mathcal{I}, [\![ \psi ]\!]_{\mathcal{I}} = 1$
\end{definition}
\begin{corollary}
  \leavevmode
  \begin{enumerate}
  \item Теория $\mathcal{T}$ выполнима (SAT), если $\exists \ \mathcal{M}$ --- модель теории $\mathcal{T}$.
  \item Теория $\mathcal{T}$ невыполнима (UNSAT), если $\nexists$ модели.
  \end{enumerate}
\end{corollary}  
\begin{remark}[Упрощение обозначения логического следствия]
  \leavevmode
  \begin{enumerate}
  \item $\{p, q\} \models r \iff p, q \models r$
  \item $T \cup \phi \models \psi \iff T , \phi \models \psi$
  \item $\varnothing \models \phi \iff \ \models \phi$
  \end{enumerate}
  Здесь $\Leftrightarrow$ --- равносильность утверждений
\end{remark}

\begin{definition}
  Пусть $A$ и $T$ --- теории.
  $A$ --- аксиома теории $T$,
  если $\forall \phi \in A, T \models \phi \ $  и  $ \ \forall \phi \in T, A \models \phi$
\end{definition}
\begin{definition}
 $T$  --- конечно аксиоматизируема, если $\exists A : A$ --- конечная аксиома теории $T$
\end{definition}

\begin{theorem}[О логическом следствии]
  Пусть $\Gamma$ --- множество формул, $\phi$ и $\psi$ --- формулы. Тогда
  \[ \Gamma, \phi \models \psi \hspace{1cm} \text{тогда и только тогда, когда} \hspace{1cm} \Gamma \models \phi \to \psi\]
\end{theorem}
\begin{proof}
  \textbf{Необходимость}:

  Пусть $\Gamma, \phi \models \psi$ и $\mathcal{I}$ --- модель $\Gamma$

  \underline{Случай 1}: $[\![ \phi ]\!]_{\mathcal{I}} = 0$.
  Очевидно, т.к. из лжи следует всё что угодно
  
  \underline{Случай 2}: $[\![ \phi ]\!]_{\mathcal{I}} = 1$.
  
  $[\![ \phi ]\!]_{\mathcal{I}} = 1 \Rightarrow [\![ \psi ]\!]_{\mathcal{I}} = 1 \Rightarrow [\![ \phi \to \psi ]\!]_{\mathcal{I}} = 1$
  
  \textbf{Достаточность}:

  Пусть $\Gamma \models \phi \to \psi$ и $\mathcal{I}$ --- модель $\Gamma$

  Знаем, что $[\![ \phi \to \psi ]\!]_{\mathcal{I}} = 1$

  Если выполнено $[\![ \phi ]\!]_{\mathcal{I}} = 1$, тогда выполнено $[\![ \psi ]\!]_{\mathcal{I}} = 1$, что удовлетворяет $\Gamma, \phi \models \psi$ по определению
  
\end{proof}

\begin{definition}
  $\phi \sim \psi \iff \phi \models \psi \ $ и $\ \psi \models \phi$
\end{definition}
\begin{example}
  \leavevmode
  \begin{itemize}
  \item $\phi \to \psi \sim \lnot \psi \to \lnot \phi$
  \item $\lnot \lnot \phi \sim \phi$
  \item $\lnot (\phi \land \psi) \sim \lnot \phi \lor \lnot \psi$
  \end{itemize}
\end{example}

\section{Классическое исчисление}
Исчисление --- это некая система над символами, дающая нам возможность дедуцировать слова из уже существующих слов, используя так называемые правила вывода.

\begin{example}
  Тут очень подошёл бы пример с SAT солвером
\end{example}

\begin{definition}
  Исчисление $\mathcal{C}$ задаётся:
  \begin{enumerate}
  \item Языком 
  \item Аксиомами
  \item Правилами вывода
  \end{enumerate}
  Теперь по-подробнее:
  
  Пусть $L$ --- язык исчисления $\mathcal{C}$ над алфавитом $\Sigma$
  \begin{enumerate}
  \item Множество аксиом $Ax$--- это подмножество $L$
  \item
    Правило вывода $\mathcal{R}$ --- это $(n + 1)$-арное отношение, где $n \in \mathbb{N}_+$

    Если $(\alpha_1, \alpha_2, ..., \alpha_n, \alpha) \in \mathcal{R}$, то говорят $\alpha$ получился из $\alpha_1, \alpha_2, ..., \alpha_n$ применением правилом вывода $\mathcal{R}$ или $\alpha$ --- результат применения $\mathcal{R}$ из $\alpha_1, \alpha_2, ..., \alpha_n$
    
  \end{enumerate}
\end{definition}

Мы будем дальше говорить, что гипотеза $\Gamma$ --- это подмножество $L$

\begin{definition}
  Вывод из гипотезы $\Gamma$ является последовательность слов $\alpha_1, \alpha_2, ..., \alpha_n$ c запятой-разделителем в которой каждое слово $\alpha_i$...
  \begin{enumerate}
  \item Либо аксиома
  \item Либо гипотеза
  \item Либо получается по некоторому правилу выводу из прошлых выведенных слов $\alpha_{i_1}, \alpha_{i_2}, ..., \alpha_{i_k}$, где $i_1, ..., i_k $ меньше $i$
  \end{enumerate}

\end{definition}

\begin{notation}
  $\Gamma \vdash \phi$ --- значит $\phi$ выводится из гипотез, принадлежащих $\Gamma$

  Ещё $\varnothing \vdash \phi \iff \vdash \phi$
\end{notation}

\begin{definition}
  Исчисления высказываний Гильбертского типа --- это исчисление, которое строится следующим образом:
  \begin{enumerate}
  \item Язык --- это язык высказываний
  \item Аксиомы: %% (Дальше пояснения в скобках --- не просто мнемоника, а соответствие Чёрча-Ховарда (TODO: Проверить на достоверность))

    Пусть $\alpha, \beta, \gamma$ --- well formed formulas

    \begin{enumerate}
    \item $\alpha \rightarrow \beta \rightarrow \alpha$ (Как комбинатор K из SKI)
    \item $(\alpha \rightarrow \beta \rightarrow \gamma) \rightarrow (\alpha \rightarrow \beta) \rightarrow (\alpha \rightarrow \gamma)$ (Как комбинатор S из SKI)
    \item $\alpha \land \beta \rightarrow \alpha$ (Как функция fst из OCaml)
    \item $\alpha \land \beta \rightarrow \beta$ (Как функция snd из OCaml)
    \item $\alpha \rightarrow \alpha \lor \beta$ (Left)
    \item $\beta \rightarrow \alpha \lor \beta$ (Right)
    \item $\alpha \rightarrow \beta \rightarrow \alpha \land \beta$ (Создание пары)
    \item $(\alpha \rightarrow \gamma) \rightarrow (\beta \rightarrow \gamma) \rightarrow (\alpha \lor \beta \rightarrow \gamma)$ (Pattern matching)
    \item $\alpha \rightarrow \top$
    \item $\bot \rightarrow \alpha$
    \item $(\alpha \rightarrow \bot) \rightarrow \lnot \alpha$
    \item $\lnot \alpha \rightarrow (\alpha \rightarrow \bot)$
    \item $(\alpha \rightarrow \beta) \rightarrow (\alpha \rightarrow \lnot \beta) \rightarrow \lnot \alpha$
    \item $\lnot \lnot \alpha \rightarrow \alpha$ (Снятие двойного отрицания)
    \end{enumerate}
  \item Правило вывода только одно --- modus ponens, а именно
  \begin{prooftree}
    \AxiomC{$\alpha, \ \alpha \rightarrow \beta$}
    \UnaryInfC{$\beta$}
  \end{prooftree}
  \end{enumerate}
\end{definition}

\begin{example}
  Зная, что первые две аксиомы исчисление Гильбертского типа --- типы функций $K$ и $S$ из комбинаторной логики SKI соответсвенно, мы можем дедуцировать $\alpha \rightarrow \alpha$

  $I = (SK)K$
\end{example}

\end{document}
